\subsection{The finite separation consequence and its exact audit scope}

\begin{result}{E.195}{Proved}
{A finite five-point level excludes the genuine parity-one locus}
Let
\[
 \mathcal X=\Delta^{-1}(1)\subset\widehat{\mathrm{GT}}_0,
 \qquad
 \mathcal P:\widehat{\mathrm{GT}}_0
                      \longrightarrow\widehat\Gamma_{0,5}
\tag{E.195.1}
\]
be the literal pentagon evaluation using the standard
insertion maps.  There is an open normal subgroup
\(V\triangleleft\widehat\Gamma_{0,5}\) such that
\[
          \{F\in\mathcal X:\mathcal P(F)\in V\}
                            =\varnothing .
\tag{E.195.2}
\]
One may simultaneously retain a finite characteristic
source quotient on which the historical character
\(\Delta\) is determined.  This proves the complete
existence assertion of historical GT A.4.1.
It specifies neither a numerical size bound nor
the exact compiler quotient of historical A.4.57,
and it does not assert an executed census of the
genuine source image.
\end{result}

\begin{proof}
The gentle group is a closed subgroup of the
continuous automorphism group of the finitely
generated free profinite group, with its stated
multiplier, peripheral and cusp normalization.
That automorphism group is profinite.  Indeed
continuous endomorphisms form a closed subset
of the finite power of the group recording
generator images; pairs of inverse endomorphisms
form a compact closed subset of its square.
Projecting those pairs gives the compact
automorphism group, with its finite
characteristic-quotient topology.
The gentle conditions are closed word and
compact conjugacy conditions, so the gentle
group and \(\mathcal X\) are compact.

Every insertion into the full five-point group
is continuous, and the pentagon is a finite
product of these evaluations.  Thus
\(\mathcal P(\mathcal X)\) is compact and closed.
By definition
\[
                \mathcal P^{-1}(1)
                         =\widehat{\mathrm{GT}}.
\tag{E.195.3}
\]
E.191 proves that this fibre is disjoint from
\(\mathcal X\).  The complement of
\(\mathcal P(\mathcal X)\) is therefore an
identity neighborhood in the profinite target;
it contains an open normal subgroup \(V\),
proving E.195.2.

The same source topology and continuity of
\(\Delta\) give a finite characteristic free-group
quotient whose congruence kernel, intersected
with the gentle group, lies in \(\ker\Delta\).
Consequently two genuine gentle automorphisms
with the same source record have the same
character.  Joining that record with the
finite target in E.195.2 gives the stated
common finite separation.

Conversely, any V satisfying E.195.2 implies
the vanishing in E.191, because an actual
pentagon-zero element maps to the identity
in every target quotient.  This converse
identifies the exact compactness implication.
It does not identify a list of formal finite
solutions with the image of the actual
gentle group, or replace a fixed source
locus by a varying family of lifts.
\end{proof}

\paragraph{Historical clauses that remain separate.}
E.191 proves the full vanishing clause of
historical A.7.  E.97 independently supplies
a genuine gentle element of character one,
so also proves
\[
 \widehat{\mathrm{GT}}\ \subseteq\ker\Delta
                    \ \subsetneq\widehat{\mathrm{GT}}_0 .
\]
The additional strict inclusion
\(\widehat{\mathrm{GT}}\subsetneq\ker\Delta\)
is independently proved in E.197--E.198.
Those results supply the actual integral and
full profinite gentle lift, followed by a
finite pentagon detector; they do not use
the historical B.2.1--B.2.3 register as a
substitute for that lifting proof.
The explicit quotient and numerical
refinement bound in A.4.57 require their
own finite construction.

The finite-versus-uniform distinction in
historical C.5 now has an established GT
premise: the actual parity-one gentle locus
is excluded at some single finite five-point
level.  On the section side E.20 retains
any prescribed finite geometric quotient
while varying the ordinary completed pair.
These assertions concern different compact
domains and different quantifiers; neither
changes the unresolved all-depth question
for one ordinary completed pair.
